\section{Kの定義}

\subsection{Kの位置付け}

$K$ は、Solver集合上に配置された公理について、その配置を基礎として公理間の関係および分離の構造を記述し、それを数値化するための量として導入する。

ここで重要なのは、「配置」「関係」「分離」をそれぞれ独立した対象として扱うのではなく、これらを一続きの構造として捉えることである。

すなわち、まずSolver集合上に各公理がどのように配置されているかを定め、その配置同士を比較することによって公理間の関係が得られ、その関係の一つの形として分離が現れる。

したがって、Kの基本的な構造は、

\[
\text{配置}
\longrightarrow
\text{配置間の関係}
\longrightarrow
\text{分離}
\longrightarrow
\text{数値化}
\]

という順序によって構成される。

このときKが示す対象は「Solver集合上に構成された公理の配置構造」であり、そこに対して最終的に数値を与える。

\subsection{Solver構造を基礎とした公理配置}

ECAにおいて、各Solver $z\in\mathcal{Z}$ には公理構造 $\mathcal{A}_z$ が対応する。

\[
z\in\mathcal{Z}
\longmapsto
\mathcal{A}_z
\]

ここで $\mathcal{A}_z$ はSolver $z$ に含まれる公理要素によって構成される。

一方、Solverには推論規則および構成制約を表す $\mathcal{R}_z$ も存在するが、Kの定義では直接の対象としない。
これは $\mathcal{R}_z$ が「推論規則および構成制約」であるのに対して、Kが扱うのは $\mathcal{A}_z$ 、すなわち「Solver $z$ における公理構造」であるためである。
より具体的には「Kが記述する対象がSolver集合上における公理の配置」だからである。

推論規則や構成制約そのものをKに直接含めた場合、Kが記述する対象が公理配置からSolverの構成全体へ拡張されるためであり、ゆえに両者を区別する必要がある。

\subsection{公理のSolver集合上での配置}

公理 $a$ に対して、その公理を含むSolverの集合を

\[
\mathcal{Z}_a = \{z\in\mathcal{Z}\mid a\in\mathcal{A}_z\}
\]

と定義する。

$\mathcal{Z}_a$ は、公理 $a$ がSolver集合 $\mathcal{Z}$ のどの部分に配置されているかを表す。

したがって、

\[
a
\longmapsto
\mathcal{Z}_a
\]

によって、公理をSolver集合上の配置へ対応させることができる。

この対応を明示的な写像として表せば、

\[
\Phi_{\mathcal{Z}}:\mathcal{A}\longrightarrow\mathcal{P}(\mathcal{Z})
\]

\[
\Phi_{\mathcal{Z}}(a) = \mathcal{Z}_a = \{z\in\mathcal{Z}\mid a\in\mathcal{A}_z\}
\]

となる。

ここで $\mathcal{A}$ は対象となる公理全体の集合である。

$\Phi_{\mathcal{Z}}$ はKそのものではない。$\Phi_{\mathcal{Z}}$ は、公理をそのSolver集合上の配置へ対応させるための構造的な写像であり、Kはその配置構造を最終的に数値化する量として位置付けられる。

\subsection{複数公理の同時配置}

複数の公理を同時に含むSolverを考えるため、$A\subseteq\mathcal{A}$ に対して

\[
\mathcal{Z}_A = \bigcap_{a\in A}\mathcal{Z}_a
\]

を定義する。

$\mathcal{Z}_a$ の定義から、

\[
\mathcal{Z}_A =\{z\in\mathcal{Z}\mid A\subseteq\mathcal{A}_z\}
\]

が得られる。

したがって、$\mathcal{Z}_A$ は、$A$ に含まれるすべての公理を同時に含むSolverの集合である。

特に、$a,b\in\mathcal{A}$ に対しては、

\[
\mathcal{Z}_{\{a,b\}} = \mathcal{Z}_a\cap\mathcal{Z}_b
\]

となる。

この構造によって、公理同士の共通配置をSolver集合上で直接表現できる。

\subsection{公理間の関係}

公理 $a,b\in\mathcal{A}$ に対して、

\[
\mathcal{Z}_a
\qquad\text{および}\qquad
\mathcal{Z}_b
\]

を比較することで、公理間の配置関係を得ることができる。

共通して配置されるSolverは、

\[
\mathcal{Z}_a\cap\mathcal{Z}_b
\]

によって与えられる。

一方、$a$ のみに対応する配置は

\[
\mathcal{Z}_a\setminus\mathcal{Z}_b
\]

であり、$b$ のみに対応する配置は

\[
\mathcal{Z}_b\setminus\mathcal{Z}_a
\]

である。

さらに、両者の配置の差をまとめて表す場合には対称差

\[
\mathcal{Z}_a\triangle\mathcal{Z}_b = (\mathcal{Z}_a\setminus\mathcal{Z}_b)\cup(\mathcal{Z}_b\setminus\mathcal{Z}_a)
\]

を用いることができる。

これらは、公理の配置から導かれる集合構造である。

また、

\[
\mathcal{Z}_a\subseteq\mathcal{Z}_b
\]

であれば、

\[
\forall z\in\mathcal{Z},\quad
a\in\mathcal{A}_z
\Longrightarrow
b\in\mathcal{A}_z
\]

が成り立つ。

このように、包含関係も公理間の配置関係として自然に得られる。

\subsection{分離の位置付け}

公理 $a,b$ の配置が完全に重ならない場合、

\[
\mathcal{Z}_a\cap\mathcal{Z}_b = \varnothing
\]

となる。

これは、公理 $a$ と公理 $b$ がSolver集合上で完全に分離されていることを表す。

したがって、分離は配置とは独立した第三の対象として導入されるのではなく、

\[
\mathcal{Z}_a
\quad\text{と}\quad
\mathcal{Z}_b
\]

という二つの配置を比較した結果として現れる。

すなわち、

\[
\text{配置}
\longrightarrow
\text{配置の比較}
\longrightarrow
\text{関係}
\longrightarrow
\text{分離}
\]

という構造になる。

特に、

\[
\mathcal{Z}_a\cap\mathcal{Z}_b=\varnothing
\]

は、この関係構造の中で分離が最大限に現れた特殊な場合として位置付けられる。

\subsection{配置の大きさと配置そのものの区別}

$\mathcal{Z}_a$ の要素数

\[
|\mathcal{Z}_a|
\]

は、公理 $a$ を含むSolverの数を表す。

したがって、有限なSolver集合において

\[
|\mathcal{Z}_a|
\]

を調べることで、公理 $a$ がどの程度の数のSolverに現れるかを知ることができる。

しかし、

\[
|\mathcal{Z}_a| = |\mathcal{Z}_b|
\]

であっても、

\[
\mathcal{Z}_a = \mathcal{Z}_b
\]

とは限らない。

例えば、

\[
\mathcal{Z}_a=\{z_1,z_2\},
\qquad
\mathcal{Z}_b=\{z_3,z_4\}
\]

であれば、

\[
|\mathcal{Z}_a| = |\mathcal{Z}_b| = 2
\]

である一方、

\[
\mathcal{Z}_a\cap\mathcal{Z}_b = \varnothing
\]

である。

したがって、配置の「大きさ」と配置の「構造」は区別する必要がある。

この区別は、Kを数値化する際にも重要である。Kの単一の数値が配置構造のすべてを表現するとは限らず、必要に応じて $\mathcal{Z}_a$ およびその間の集合関係を基礎構造として保持したうえで、数値評価を行う必要がある。

\subsection{Kの基本的定義}

以上を踏まえると、Kは次のように定義できる。

\[
\boxed{
K\text{ は、Solver集合上に配置された公理について、その配置を基礎として公理間の関係および分離の構造を記述し、それを数値化する量である。}
}
\]

この定義では、「配置」「関係」「分離」は独立した三つの選択肢ではない。

まず、

\[
\mathcal{Z}_a = \{z\in\mathcal{Z}\mid a\in\mathcal{A}_z\}
\]

によって公理の配置を定める。

次に、

\[
\mathcal{Z}_a\cap\mathcal{Z}_b,
\qquad
\mathcal{Z}_a\setminus\mathcal{Z}_b,
\qquad
\mathcal{Z}_b\setminus\mathcal{Z}_a,
\qquad
\mathcal{Z}_a\triangle\mathcal{Z}_b
\]

などによって配置間の関係を得る。

そして、

\[
\mathcal{Z}_a\cap\mathcal{Z}_b=\varnothing
\]

という関係の特殊な場合として分離が得られる。

最後に、この配置構造を数値評価することでKの値が得られる。

したがって、Kの構造は

\[
\boxed{
\mathcal{Z}_a
\longrightarrow
\{\mathcal{Z}_a\mid a\in\mathcal{A}\}
\longrightarrow
\text{配置間の関係}
\longrightarrow
\text{分離}
\longrightarrow
K
}
\]

として整理できる。

\subsection{KとB,Iとの関係}

KはSolver集合上の公理配置を数値化する量であるが、Kだけで公理体系全体の構造を表現するものではない。

ECAでは、B、I、Kはそれぞれ異なる側面を扱う。

Bは公理数や公理選択空間に関する量であり、Solver側の公理構成の規模を扱う。

Iは公理体系間の構造整合性を扱う。

これに対してKは、

\[
\text{公理}
\longrightarrow
\text{Solver集合上の配置}
\]

という方向から、公理がSolver集合上にどのように配置されているかを扱う。

したがって、

\[
B,\ I,\ K
\]

は同一の情報を別々の方法で表現するものではなく、それぞれ異なる構造を数値化する。

この区別によって、K単独では表現しきれない情報についても、ECAの数値評価

\[
T_{\mathrm{overlay}} = \frac{1}{N}\sum_{i=1}^{N}B_iI_iK_i
\]

においてBおよびIと組み合わせて扱うことができる。

したがって、Kに対して「K単独の数値だけで公理配置の全構造を完全に復元できなければならない」という条件を課す必要はない。

\subsection{定義から数値化への接続}

集合論におけるKの定義段階では、

\[
\mathcal{Z}_a
\]

および

\[
\{\mathcal{Z}_a\mid a\in\mathcal{A}\}
\]

によってSolver集合上の公理配置を構成する。

その後、配置構造に対して数値評価を導入する。

したがって、Kの定義とKの数値化は次のように区別される。

\[
\boxed{
\text{集合構造}
\longrightarrow
\text{配置関係}
\longrightarrow
\text{数値評価}
}
\]

この区別により、KをECAのSolver集合構造から自然に数値化することができる。